\documentclass[a4paper,review,12pt]{elsarticle}

\usepackage{booktabs}
\usepackage{amsmath}
\usepackage{graphicx}
\usepackage{hyperref}

\journal{Image and Vision Computing}

\begin{document}

\begin{frontmatter}

\title{PACE-Seg: Platform-Aware Calibrated Elastic Semantic Segmentation for Reliable Edge Vision}

\author{D\~ung Nguy\~ên}
\affiliation{organization={}, % TODO: institution/affiliation
             addressline={},
             city={},
             country={}}

\begin{abstract}
% TODO: abstract cannot be honestly written until Introduction/Related work
% (blocked on the systematic literature review, RESEARCH_PLAN.md \S2) and the
% Ablations/Conclusion sections exist. Draft only once those are in place --
% an abstract written against an incomplete paper tends to overclaim.
TODO.
\end{abstract}

\begin{keyword}
semantic segmentation \sep edge deployment \sep elastic neural networks \sep quantization-aware training \sep calibrated uncertainty \sep hardware-aware neural architecture search
\end{keyword}

\end{frontmatter}

\section{Introduction}
% TODO: blocked on the systematic literature review (RESEARCH_PLAN.md \S2) --
% any novelty claim stronger than a safe "we study the underexplored
% intersection of..." framing needs the Scopus/Web of Science search that
% section requires, not yet done. See docs/MANUSCRIPT_SKELETON.md's
% Introduction section for the safe interim framing and the four gaps
% (G1-G4) this work addresses, kept out of this .tex file until the review
% is done and the framing can be written as real prose rather than restated
% from the skeleton verbatim.
TODO.

\section{Related work}
% TODO: blocked on the same literature review as Introduction. Needs: (1)
% lightweight/real-time segmentation and the FLOPs-vs-latency disconnect;
% (2) compiler-portability across TensorRT, DLA and Hailo-style dataflow
% NPUs; (3) adverse-condition segmentation/uncertainty benchmarks (ACDC);
% (4) dynamic/once-for-all supernet and routing literature. See
% docs/RESEARCH_PLAN.md \S2 and docs/SOURCE_RESEARCH_GAP_2026.md.
TODO.

\section{Method}
\label{sec:method}

\subsection{Compiler-safe elastic supernet}
\label{sec:supernet}

We train a single shared segmentation network that executes at four discrete
elasticity levels --- \emph{tiny}, \emph{small}, \emph{medium}, \emph{large}
--- obtained by varying channel width multiplier, block count, and input
resolution. Every operator in the network (convolution--batch-norm--activation,
depthwise/pointwise convolution, pooling, elementwise add/concatenate,
static-factor bilinear resize) is drawn from the intersection of operators we
verified compile and execute correctly on all three target compiler backends
(Section~\ref{sec:compiler-validation}), rather than from an unconstrained
search space later found to be unsupported on one target. All four levels are
trained jointly with a sandwich-rule schedule (training the smallest, largest,
and a randomly sampled subset of middle levels at each step) and in-place
distillation from the largest level's output to the smaller levels, with a
boundary-aware auxiliary loss term that up-weights pixels near class
boundaries --- the region where small classes are disproportionately lost at
low capacity.

Training uses real Cityscapes (2{,}975 train / 500 val) and ACDC (1{,}600
train / 406 val, four adverse conditions: fog, night, rain, snow) images at a
batch size and learning rate schedule held fixed across every experiment in
this paper. We additionally apply random scale-crop-pad, horizontal flip, and
color jitter augmentation to the training stream; Section~\ref{sec:aug}
quantifies this choice's effect across five different segmentation
architectures rather than assuming it helps uniformly.
% source: README.md Phase 4; reports/augmentation_effect_all_5_baselines_20260916.md

\subsection{Hardware-in-the-loop Pareto subnet selection}
\label{sec:pareto}

Rather than selecting among the four trained levels by parameter count or
FLOPs, we select by measured deployment cost. For each of four deployment
targets we benchmarked --- a Hailo-8 M.2 accelerator (Raspberry Pi~5 host) and
three NVIDIA Jetson boards spanning three tiers (Xavier NX, AGX Xavier, Orin
Nano; TensorRT, FP16) --- we measure end-to-end and kernel-only latency for
all four levels, pooling three independent runs per level per the protocol in
Section~\ref{sec:latency-protocol}, and build a per-target Pareto frontier
over the (latency, mIoU) plane. On every target, all four levels are
Pareto-optimal (no level is simultaneously slower and less accurate than
another); the level that is optimal under a fixed 10~ms latency budget differs
by target --- \emph{small} on the Hailo-8, \emph{small} on Xavier NX,
\emph{medium} on Orin Nano, \emph{large} on AGX Xavier, whose TensorRT GPU
path is the fastest of the four by a wide margin. Four real devices producing
three different choices at one identical budget is concrete evidence that
hardware-aware subnet selection is not interchangeable with a single,
device-agnostic choice --- the premise our approach is built on --- though we
have not yet quantified this against a FLOPs-aware selection baseline
(Section~\ref{sec:limitations} discusses this gap directly rather than
treating the qualitative result as sufficient on its own).
% source: reports/pareto_search_v1_20260912.md; reports/edge/E2_E5_tensorrt_benchmark_all_levels_20260916.md

\subsection{Quantization-aware training}
\label{sec:qat}

We convert every convolution --- both the plain convolutions in our
comparison baselines and the width-sliced slimmable convolutions specific to
the shared supernet --- to a fake-quantized form: dynamic, per-tensor,
symmetric INT8 quantization of each layer's weight and input activation,
applied with a straight-through gradient estimator so training proceeds
unmodified through the quantization op. For the supernet's slimmable
convolutions, the quantization range is computed from each elasticity level's
own \emph{active} channel slice at every forward call, so every level
receives an appropriately scaled range without any additional bookkeeping. We
deliberately use a dynamic, uncalibrated range rather than a range fit from a
held-out calibration set spanning all five visual conditions
(Section~\ref{sec:supernet}) --- a documented simplification that should be
read as an upper bound on achievable accuracy, not a final number, until a
calibrated variant is built.

We fine-tune each INT8 model for 10{,}000 steps from a converged FP32
checkpoint at a reduced learning rate ($3\times10^{-5}$), rather than training
INT8 from random initialization. Across two representative fixed-budget
architectures --- Fast-SCNN (1.14~M parameters) and DDRNet-23-slim (5.22~M
parameters) --- and two random seeds each, INT8 fine-tuning loses between 0.11
and 1.25 mIoU points relative to the FP32 checkpoint it was fine-tuned from,
across Cityscapes and all four ACDC conditions (worst case: Fast-SCNN seed~0
on rain, $-1.25$ points). This is within the go/no-go accuracy budget we fixed
in advance ($\leq 1.0$--$1.5$ mIoU points; Section~\ref{sec:go-no-go}),
confirmed across two architectures and two seeds rather than a single run ---
for independent, fixed-architecture baselines. Applying the same recipe to the
shared supernet's \emph{large} level itself, across the two seeds used in
Section~\ref{sec:go-no-go}'s comparison, is \textbf{more mixed}: one seed
stays within budget on every split (worst case $-1.39$ points), while the
other exceeds it on two of five splits (Cityscapes $-1.60$, ACDC/night
$-1.66$) --- worse than any independent baseline result. We do not yet have
enough seeds to say whether this is a real, reproducible effect of quantizing
weights the sandwich rule shares across elasticity levels, or seed noise, and
report it as an open question rather than folding it into the ``confirmed''
claim above.
% source: reports/qat_v1_20260913.md (2026-09-14/15/16 updates) -- do not
% average away the seed0 overshoot when citing this result

\subsection{Compiled static engines and calibrated routing}
\label{sec:router}

Each trained elasticity level is exported and compiled independently ---
never as a single graph with runtime branching --- to a TensorRT engine, a
DLA-targeted engine, and a Hailo HEF (via ONNX and the Hailo Dataflow
Compiler); a deployed system picks one compiled engine per inference window
rather than executing dynamic control flow inside a compiled graph. This
partition into independently-compiled static engines is what makes
hardware-aware routing possible without a graph any of the three compiler
toolchains would reject.

The routing decision is made by a lightweight, calibrated risk probe rather
than a fixed policy. We run one cheap forward pass --- typically at the
\emph{tiny} level --- and take the mean per-pixel softmax entropy of its
output as a raw, uncalibrated risk score; we then map this score to an
expected pixel-error probability using a monotonic calibrator
(quantile-binned with a cumulative-max correction) fit on a held-out half of
each validation split, never on the images a policy is later evaluated
against. At inference time, a target risk budget determines how many
elasticity levels to escalate above the cheapest available level.

We evaluate this policy end to end: for the held-out half of each validation
split, we run every candidate level's forward pass, route each image under
each candidate policy using only its cheap-pass risk score, and measure the
resulting \emph{achieved} mIoU and the real measured latency
(Section~\ref{sec:pareto}) the policy actually spends. Calibrated routing
achieves higher mIoU than routing on raw, uncalibrated entropy in every one of
seven tested conditions --- Cityscapes under two risk-budget settings, two
deployment targets (Hailo-8 and AGX Xavier), and all four ACDC adverse
conditions --- with margins of $+0.0013$ to $+0.0183$ mIoU. The efficiency
picture is more mixed: because calibrated routing also spends slightly more
latency than entropy-only routing in every condition, the
mIoU-gained-per-millisecond-spent comparison favors calibrated routing in two
of the seven conditions and entropy-only routing in the other five. We report
both numbers rather than only the mIoU comparison that favors our method.

A structural property of the current policy is worth stating plainly:
escalation decisions depend only on each candidate level's latency
\emph{rank}, not its magnitude, so a calibrator fit once transfers unchanged
across deployment targets --- confirmed empirically (Hailo-8 and AGX Xavier
produce numerically identical achieved mIoU on the same validation split,
despite a $>\!4\times$ difference in absolute per-level latency). This is a
useful robustness property --- no per-device recalibration is required --- but
it also means the current policy cannot yet exploit a target-specific cost
\emph{ratio} between levels, which we flag as a direction for a
latency-value-aware successor policy rather than treating as resolved.
% source: reports/router_v1_20260914.md

\section{Deployment protocol}
\label{sec:deployment}

\subsection{Compiler validation}
\label{sec:compiler-validation}

We validate each of the three target compiler toolchains independently before
relying on any of them for a result: TensorRT (GPU and DLA paths, via
\texttt{trtexec}), and the Hailo Dataflow Compiler (ONNX~$\to$~HAR~$\to$~HEF).
TensorRT compiles and executes all four elasticity levels in both FP16 and
INT8 precision on two Jetson devices (8/8 configurations pass); the Hailo
toolchain compiles all four levels to a working HEF, and all four ran
successfully on physical Hailo-8 hardware. Xavier's DLA path compiles only the
encoder half of the network at every elasticity level --- the decoder always
falls back to the GPU. We traced this to a hard, documented hardware limit (a
DLA core executes at most 16 subgraphs, a budget the encoder alone exhausts),
not an unsupported operator, ruling out a fix by further restricting our
operator set. Because this contingency was anticipated in our go/no-go
criteria (excessive DLA fallback), we demote DLA to a secondary,
encoder-only ablation and report TensorRT GPU and Hailo HEF as the two primary
deployment backends for every headline result in this paper.
% source: README.md Phase 2; reports/edge/E3_compiler_smoke_test_20260908.md; reports/edge/E1_hailo_dfc_compile_20260909.md

\subsection{Latency measurement protocol}
\label{sec:latency-protocol}

Latency is measured with a protocol adapted from MLPerf Power's methodology
(without claiming formal compliance): batch size~1; identical resolution and
pre/post-processing held fixed within a configuration; a warm-up period before
timed measurement; three independent runs per (device, backend, precision,
elasticity level) configuration, pooled into one bootstrap 95\% confidence
interval rather than reported as three separate point estimates; kernel-only
and end-to-end latency reported separately, never compared across device
families where the two are not measuring the same boundary. We report the
complete measured cross-backend latency table for all four benchmarked
targets across all four elasticity levels in Section~\ref{sec:latency-results}.

\textbf{Energy is not yet reported.} The Hailo-8 M.2 module used in this work
exposes no on-board power or current telemetry (\texttt{--measure-power} and
\texttt{--measure-current} are both unsupported on this hardware), and we do
not yet have access to an external calibrated power meter for any target ---
a material gap against our own protocol, disclosed rather than worked around
with an unaudited on-chip estimate.
% source: README.md Phase 3; reports/edge/E1_hailo_benchmark_protocol_20260910.md; reports/edge/E3_tensorrt_benchmark_all_levels_20260911.md; RESEARCH_PLAN.md \S9

\section{Experiments}
\label{sec:experiments}

\subsection{Vision quality is monotonic in elasticity level}
\label{sec:monotonic}

We first confirm the basic premise the elastic supernet depends on: that its
four trained levels form a genuine accuracy ladder, not four
arbitrarily-ordered configurations. At 100{,}000 training steps, mIoU is
monotonic in model size at every one of five evaluation splits
(Table~\ref{tab:monotonic}), and reproduces within 0.01--0.03 mIoU on an
independently seeded run.

\begin{table}[h]
\centering
\caption{Supernet mIoU by elasticity level (seed 0, 100k steps).}
\label{tab:monotonic}
\begin{tabular}{lrrrrr}
\toprule
Level & Cityscapes & ACDC/fog & ACDC/night & ACDC/rain & ACDC/snow \\
\midrule
tiny   & 0.2986 & 0.3154 & 0.1944 & 0.2901 & 0.2610 \\
small  & 0.3564 & 0.3658 & 0.2511 & 0.3680 & 0.3304 \\
medium & 0.4120 & 0.4360 & 0.2898 & 0.4003 & 0.3981 \\
large  & 0.4712 & 0.4919 & 0.3312 & 0.4517 & 0.4492 \\
\bottomrule
\end{tabular}
\end{table}

ACDC/night is the hardest condition throughout, as expected of the
lowest-light, lowest-contrast adverse condition. That ACDC/fog scores
\emph{above} clean Cityscapes at every level is a real effect we traced to
class composition rather than an artifact of unweighted macro-averaging:
several large structural classes (wall, pole, traffic light, sky) score
markedly higher in fog's visually simpler scenes, while dynamic road-user
classes degrade as expected under fog (at the \emph{large} level: person
$0.471\to0.280$, bicycle $0.472\to0.204$, car $0.828\to0.709$). We report
per-class breakdowns for exactly these dynamic classes alongside every
aggregate mIoU number in this paper for this reason --- an aggregate number
alone would mask a real safety-relevant degradation behind an unrelated
class-mix effect.
% source: reports/first_full_supernet_run_100k_20260910.md

\subsection{Shared supernet vs.\ independent per-budget training}
\label{sec:go-no-go}

RQ2 asks whether a single shared supernet can match independently trained,
same-parameter-budget models without the cost of training and maintaining one
model per target --- not whether it can beat them. We compare the supernet's
\emph{large} level (1.02~M parameters) against Fast-SCNN (1.05~M parameters,
the one comparison baseline at a matched budget), both trained with identical
augmentation, over three independent seeds per side, per our go/no-go
criterion that headline accuracy numbers require three training seeds.

\begin{table}[h]
\centering
\caption{Supernet-large vs.\ Fast-SCNN, augmented, 3-seed average.}
\label{tab:go-no-go}
\begin{tabular}{lrrr}
\toprule
Split & Fast-SCNN (aug) & Supernet-large (aug) & Gap \\
\midrule
Cityscapes  & 0.5228 & 0.5338 & $+0.0110$ \\
ACDC/fog    & 0.5629 & 0.5650 & $+0.0021$ \\
ACDC/night  & 0.3822 & 0.3757 & $-0.0065$ \\
ACDC/rain   & 0.5050 & 0.5027 & $-0.0023$ \\
ACDC/snow   & 0.5048 & 0.5055 & $+0.0007$ \\
\bottomrule
\end{tabular}
\end{table}

The supernet wins three of five splits and loses two, with every margin
within 1.1~mIoU points --- a near-exact tie on ACDC/snow. We read this as
\textbf{true near-parity, not a win for either side}: the shared supernet
matches independently trained, same-budget training within roughly one mIoU
point across clean and all four adverse conditions, while requiring one
training run instead of four (or four times four, across our elasticity
levels). This is RQ2's actual hypothesis, confirmed rather than exceeded, and
is comfortably clear of our own no-go trigger ($>2$ mIoU loss in either
direction). We note this result required three successive corrections during
development --- an un-augmented, single-seed comparison had initially favored
the supernet by 2.6--6.4 mIoU, an artifact of Fast-SCNN's own augmentation gap
rather than a real capability difference --- and report only the final,
3-seed, augmented-vs-augmented number as citable.
% source: README.md Phase 8; reports/baseline_comparison_gap_check_20260912.md

\subsection{Data augmentation's effect is architecture-dependent, not monotonic in model size}
\label{sec:aug}

Motivated by an unexplained result on one architecture during the comparison
above, we trained all five of our required baseline architectures both with
and without our augmentation policy (random scale-crop, horizontal flip,
color jitter), spanning almost an order of magnitude in parameter count
(1.14~M--11.0~M). The effect is not a monotonic function of size
(Table~\ref{tab:aug}).

\begin{table}[h]
\centering
\caption{Augmentation effect (augmented $-$ non-augmented mIoU) by architecture.}
\label{tab:aug}
\begin{tabular}{lrrrr}
\toprule
Model & Params (M) & No-aug Cityscapes mIoU & Cityscapes $\Delta$ & ACDC-avg $\Delta$ \\
\midrule
Fast-SCNN                & 1.14  & 0.408 & $\mathbf{+0.115}$ & $\mathbf{+0.107}$ \\
BiSeNetV2                & 1.71  & 0.581 & $-0.004$ & $+0.003$ \\
SegFormer-B0              & 3.72  & 0.567 & $-0.015$ & $+0.016$ \\
DDRNet-23-slim            & 5.22  & 0.627 & $+0.000$ & $-0.005$ \\
MobileNetV3+DeepLabV3     & 11.03 & 0.577 & $\mathbf{+0.026}$ & $\mathbf{+0.045}$ \\
\bottomrule
\end{tabular}
\end{table}

The smallest model (Fast-SCNN) is, without augmentation, badly underfit ---
its 0.408 Cityscapes mIoU is far the worst of the five --- and augmentation
relieves this directly, producing the largest gain of any architecture. The
\emph{largest} model (MobileNetV3+DeepLabV3) also gains substantially, despite
showing no sign of underfitting: its non-augmented mIoU (0.577) is in fact
\emph{worse} than DDRNet-23-slim's (0.627), a model with under half its
parameters, a pattern consistent with overfitting our comparatively small
(4{,}575-image) training set --- a failure mode augmentation's regularizing
effect also relieves, through a different mechanism than Fast-SCNN's. The
three architectures between these extremes (1.7--5.2~M parameters, the three
best-fitting models without augmentation) show small, mixed, and sometimes
slightly negative effects. We read this as a bimodal relationship with how
comfortably a given architecture fits this particular dataset and training
budget, not a monotonic function of parameter count --- and caution against
citing ``augmentation helps small models more'' as a general claim without
first establishing whether a given architecture is under- or over-fit at its
evaluated budget.
% source: reports/augmentation_effect_all_5_baselines_20260916.md

\subsection{Measured latency and hardware-aware subnet selection}
\label{sec:latency-results}

\begin{table}[h]
\centering
\caption{End-to-end latency (ms) by elasticity level, mean over 3 independent runs, across all four benchmarked deployment targets.}
\label{tab:latency}
\begin{tabular}{lrrrr}
\toprule
Level & Hailo-8 & Xavier NX (TensorRT FP16) & AGX Xavier (TensorRT FP16) & Orin Nano (TensorRT FP16) \\
\midrule
tiny   & 3.714 & 2.061 & 0.912 & 1.081 \\
small  & 6.597 & 4.964 & 1.781 & 2.604 \\
medium & 20.29 & 15.10 & 4.546 & 7.727 \\
large  & 41.14 & 40.26 & 9.389 & 17.95 \\
\bottomrule
\end{tabular}
\end{table}

Joining this table with Table~\ref{tab:monotonic}'s mIoU numbers per level
(Section~\ref{sec:pareto}), all four levels lie on the Pareto frontier on
every target --- none is simultaneously slower and less accurate than
another --- and the level selected under a fixed 10~ms budget differs by
target: \emph{small} on the Hailo-8, \emph{small} on Xavier NX, \emph{large}
on AGX Xavier (whose TensorRT path is the fastest of the four, 4--4.5$\times$
the Hailo-8's per-level latency), and \emph{medium} on Orin Nano. Four real
devices producing three distinct choices at one identical latency budget is
the full device-lineup evidence our protocol calls for --- stronger than a
two-device demonstration alone, though still only a qualitative confirmation
that hardware-aware selection matters, not yet the quantitative comparison
against a FLOPs-aware baseline that would let us claim a specific margin
(Section~\ref{sec:limitations}).
% source: reports/pareto_search_v1_20260912.md; reports/edge/E2_E5_tensorrt_benchmark_all_levels_20260916.md

\section{Ablations and failure analysis}
\label{sec:ablations}
% TODO: blocked on Phase 9 (RESEARCH_PLAN.md \S10's 10 ablation axes), only
% ~10-15\% started (a handful of axes touched as side effects of other
% sections: shared-vs-independent training in \S\ref{sec:go-no-go}, entropy
% vs.\ calibrated risk in \S\ref{sec:router}). A with/without-distillation
% ablation run and a second QAT-on-supernet seed are in flight as of
% 2026-09-16 -- fill this section in once they complete and are evaluated.
TODO.

\section{Limitations}
\label{sec:limitations}
% TODO: draft once Ablations exists -- premature to enumerate limitations
% before the ablation results that would motivate several of them (e.g. the
% FLOPs-aware baseline gap, the latency-rank-only router policy, the
% single-target-pair QAT evidence) are in hand.
TODO.

\section{Conclusion}
\label{sec:conclusion}
% TODO: blocked on the above.
TODO.

\section*{Traceability note (remove before submission)}
Every quantitative claim in this draft is annotated with a \texttt{\% source:}
LaTeX comment pointing to the specific report under \texttt{reports/} in the
\texttt{imavis-edge-seg} repository it was drawn from. Before final
submission, verify no number has drifted from its source during editing, then
delete this section and all \texttt{\% source:} comments.

\end{document}
